\subsection{能量、方差与褐变}
下列数字由 \texttt{s1\_summary.json} 的 19440 行重新分组。能量仍是加热指令乘记录间隔。方差低于 0.3 表示走过一次乘 0.35 的更新；未更新的先验方差是 0.64。褐变是终点 $\alpha_b$，不是停机时刻。
\begin{longtable}{lrrrrrr}
\caption{按厚度与策略的能量分位与终点方差}\\
\toprule
分组 & $n$ & $E_{10}$ & $E_{50}$ & $E_{90}$ & 方差$<0.3$ & $\alpha_{b,50}$ \\
\midrule\endfirsthead
\toprule 分组 & $n$ & $E_{10}$ & $E_{50}$ & $E_{90}$ & 方差$<0.3$ & $\alpha_{b,50}$ \\\midrule\endhead
15\,mm, 开环 & 1620 & 900 & 900 & 900 & 0.000 & 0.0000 \\
15\,mm, 腔温比例 & 1620 & 3310 & 3310 & 3310 & 0.000 & 0.0000 \\
15\,mm, 无探针 & 1620 & 892 & 933 & 1076 & 0.000 & 0.0000 \\
15\,mm, 对偶 & 1620 & 890 & 932 & 1064 & 1.000 & 0.0000 \\
30\,mm, 开环 & 1620 & 1170 & 1170 & 1170 & 0.000 & 0.0000 \\
30\,mm, 腔温比例 & 1620 & 3310 & 3310 & 3310 & 0.000 & 0.0000 \\
30\,mm, 无探针 & 1620 & 1563 & 1655 & 1753 & 0.000 & 0.0000 \\
30\,mm, 对偶 & 1620 & 1629 & 1730 & 1836 & 1.000 & 0.0000 \\
50\,mm, 开环 & 1620 & 1530 & 1530 & 1530 & 0.000 & 0.0000 \\
50\,mm, 腔温比例 & 1620 & 3310 & 3310 & 3310 & 0.000 & 0.0000 \\
50\,mm, 无探针 & 1620 & 2901 & 3055 & 3178 & 0.000 & 0.0000 \\
50\,mm, 对偶 & 1620 & 3157 & 3250 & 3307 & 1.000 & 0.0000 \\
\bottomrule\end{longtable}
对偶在三档厚度上的方差下降比例都是 1。开环、腔温比例和无探针的比例是 0，因为它们的代码路径不写 $\hat r$。终点褐变的中位数在这十二格里都是 0：表面含水没有降到褐变式要求的 0.55 以下，或者表面温度没有同时越过 $110^\circ\mathrm{C}$。中位数为 0 不表示每一条轨迹的 $\alpha_b$ 都是 0。三个原型在 S1 的含水格子上活度都贴近饱和，所以这张表也不能比较蛋白与淀粉的表面反应。
